\chapter{Calibration and objections}
\label{ch:calibration}

\section*{What this chapter is for}

Two related skills: estimating honestly where you stand in a process, and
recognising which kind of objection you are facing. The second determines
whether the first can be improved by anything you do.

\section{Estimating a process}

A loop is a chain of conditional probabilities, and treating it as one is
useful mainly because it forces you to say where you think the risk is.

The mechanics are simple: assign a pass probability to each remaining stage,
multiply. What matters is not the product but the exercise \dash{} you cannot
assign a number to a stage without asking what it measures and what your
evidence for it is, which is the useful part.

\begin{kitnote}
Update after each stage, and write the number down before you learn the
outcome. An estimate recorded afterwards is not an estimate. This is the same
discipline as writing a prediction before running an experiment, for the same
reason.
\end{kitnote}

\section{The failure mode worth knowing in advance}

The naive model treats a passed stage as settled: the concern was raised, it
was addressed, it is behind you.

\textbf{It is not.} A later interviewer is not bound by an earlier one's
acceptance. A concern that three people found answerable can be weighed
differently by a fourth with different authority and different responsibility
\dash{} and the fourth is frequently the last one.

\judgement{So a raised concern never leaves your estimate. It gets a lower
weight after it has been accepted once, and it does not go to zero, and the
error of setting it to zero runs entirely in the direction of optimism. If you
find yourself confident going into a final round on the grounds that the hard
parts are behind you, that is the moment this applies to.}

The practical consequence is not pessimism. It is that a concern you know you
carry is worth addressing \emph{structurally} \dash{} by changing the evidence
\dash{} rather than by refining the explanation, because the explanation has to
work on every reader independently and the evidence only has to be true once.

\section{Two kinds of objection}

This is the distinction the chapter exists for.

\begin{kittable}{@{}lXX@{}}
& \textbf{Narrative objection} & \textbf{Data objection} \\
\midrule
Shape & ``I do not understand X'' & ``X is the case'' \\
Rests on & missing or mis-framed information & a fact about your record \\
Answered by & context, reframing, evidence they had not seen & different facts \\
Example & ``this stack does not look relevant'' & ``you have not done this at scale'' \\
Effort that helps & preparation & time \\
\end{kittable}

A narrative objection is answerable. The interviewer is missing something, or
has a reasonable reading that a better framing corrects, and preparation moves
it.

A data objection is not answered by explanation. Explaining a fact does not
change it, and \judgement{the specific trap is that explaining well
\emph{feels} like progress: the conversation goes smoothly, the interviewer
nods, and nothing about the underlying evidence has moved. An explanation can
carry you through several stages and then meet somebody who is weighing the
fact rather than the account of it.}

\begin{kitmethod}
\textbf{Tell them apart by asking: if they believed every word, would the
concern be gone?}

If yes, it is narrative. Prepare the framing, deliver it once, move on.

If no \dash{} if they would believe you completely and the concern would remain
\dash{} it is a data objection. Then: answer honestly and briefly, do not
over-explain, do not offer compensating concessions, and understand that the
process may end there for reasons that are about evidence rather than about
you.
\end{kitmethod}

\section{Answering either one}

\textbf{Narrative.} Name it before they have to. Give the missing context in
two sentences. Point at evidence. Stop talking \dash{} the commonest error is
continuing past the point where the objection was answered, which reintroduces
it.

\textbf{Data.} Acknowledge it plainly, without a wrapper of justification. Give
the honest reason once, if there is one. Say what you are doing about it, if
anything, in one sentence. Then stop.

\begin{kittrap}
When a data objection lands at a late stage, the instinct is to offer something
\dash{} a smaller engagement, a trial period, a lower level, a reduced
expectation. It feels like problem-solving and it is nearly always
counter-productive, because it answers a concern about commitment or capacity
with a proposal that is smaller and less committed.

If a genuinely different arrangement would suit you, propose it at the
beginning, where it is a preference. Proposing it at the end makes it a
concession, and a concession confirms the concern it was meant to dissolve.
\end{kittrap}

\begin{drill}
Write the two strongest objections to hiring you, in the interviewer's words.
For each, apply the test: if they believed everything you said, would the
concern be gone?

Then write the answer to each \dash{} five sentences for a narrative objection,
two for a data one. If both come out as data objections, the useful work is
Chapter~\ref{ch:classes} rather than another rehearsal.
\end{drill}

\section*{What this chapter does not claim}

The estimation method is not a forecasting technique and no accuracy is claimed
for it. It is a device for making your own assumptions visible.

The claim that a concern can return at a later stage is the author's judgement,
supported by the structure of a loop \dash{} later interviewers have different
authority and are not bound by earlier acceptances \dash{} rather than by any
measurement of how often it happens.
